% =========================================================
% TABLA 1: OBJETIVO ESPECÍFICO 1
% =========================================================
\begin{table}[htbp]
	\centering
	\caption{Cumplimiento del Objetivo Específico 1}
	\label{tab:cumplimiento-objetivo-1}

	\begin{tabular}{
		|>{\raggedright\arraybackslash}p{0.38\textwidth}
		|>{\raggedright\arraybackslash}p{0.56\textwidth}|}
		\hline
		\rowcolor[HTML]{1B6B93}
		\multicolumn{2}{|c|}{\color{white}\textbf{\textit{Objetivo específico 1}}}
		\\
		\hline
		\multicolumn{2}{
		|>{\centering\arraybackslash}p{\dimexpr0.94\textwidth+2\tabcolsep+\arrayrulewidth\relax}|
		}{\textit{Identificar la necesidad funcional y técnica del GRID con relación a un servicio de directorio}}
		\\
		\hline
		\textbf{\textit{Detalle de cumplimiento}}                                                                                                                                                                                                                                                                                                                                                             &
		\textbf{\textit{Observaciones}}
		\\
		\hline
		Ver capítulo \textcolor[HTML]{1B6B93}{\textbf{\ref{cap:caracterizacion-grid}}} \textbf{Caracterización del \ac{GRID}} y capítulo \textcolor[HTML]{1B6B93}{\textbf{\ref{cap:identificacion-requisitos}}} \textbf{Identificación de requisitos}. Se realizaron procesos cuali-cuantitativos para obtener información del grupo de investigación y delimitar los aspectos que debía abordar la solución. &
		\textbf{Paso 1: Análisis del entorno del grupo de investigación.} Se realizó una caracterización del grupo \ac{GRID}, incluyendo stakeholders, integrantes, áreas de trabajo, misión, visión e impacto del proyecto (capítulo \textcolor[HTML]{1B6B93}{\textbf{\ref{cap:caracterizacion-grid}}}).\par\vspace{0.4em}
		- \textbf{Entrevista} con integrante del grupo \ac{GRID} (sección \textcolor[HTML]{1B6B93}{\ref{sec:entrevista-grid}}).\par\vspace{0.6em}
		\textbf{Paso 2: Identificación de necesidades, problemas y oportunidades.}\par\vspace{0.4em}
		- Se identificaron las necesidades funcionales y técnicas relacionadas con la gestión de identidades y el servicio de directorio (sección \textcolor[HTML]{1B6B93}{\ref{subsec:identificacion-necesidades}}).\par\vspace{0.6em}
		\textbf{Paso 3: Definición de requisitos.}\par\vspace{0.4em}
		- Se definieron los requisitos funcionales R-01 a R-08 y los requisitos de seguridad asociados (sección \textcolor[HTML]{1B6B93}{\ref{sec:definicion-requisitos}}, ver \autoref{tab:requisitos}).                                                                                                                                                                                                       \\
		\hline
	\end{tabular}
\end{table}
